\documentclass[12pt, a4paper, oneside]{article}
\usepackage{amsmath, amsthm, amssymb, bm, graphicx, hyperref, mathrsfs, booktabs}
\usepackage[margin=2.5cm]{geometry}

\linespread{1.3}
\newtheorem{theorem}{Theorem}
\newtheorem{remark}[theorem]{Remark}

\newcommand{\qvec}{\mathbf{q}}
\newcommand{\rvec}{\mathbf{r}}
\newcommand{\vvec}{\mathbf{v}}
\newcommand{\qhat}{\hat{\mathbf{q}}}
\newcommand{\Nt}{N_t}
\newcommand{\Na}{N}
\newcommand{\Nw}{N_\omega}
\newcommand{\field}{S}
\newcommand{\Wmat}{\mathbf{W}}
\newcommand{\Imat}{\mathbf{I}}

\begin{document}

\thispagestyle{empty}
\vspace*{\fill}
\begin{center}
    \Huge\textbf{Density of States from the\\On-Site Time Correlation}

    \vspace{1cm}
    \Large\textbf{Algorithm and Formulas}

    \vspace{0.8cm}
    \large \texttt{corr/sqw\_dos.py}
\end{center}
\vspace*{\fill}

\newpage
\tableofcontents
\newpage

% ======================================================================
\section{Overview}
% ======================================================================

The density of states is the dynamic structure factor integrated over the
Brillouin zone,

\begin{equation}
    D(\omega) = \sum_{\qvec} S(\qvec,\omega).
\end{equation}

Computing it by evaluating $S(\qvec,\omega)$ on a $\qvec$-mesh and summing
would be wasteful, because the sum can be done analytically first. Writing the
spatial Fourier amplitude as

\begin{equation}
    s^\alpha(\qvec,t) = \frac{1}{\sqrt{\Na}}\sum_j \field_j^\alpha(t)\,
                        e^{i\qvec\cdot\rvec_j},
\end{equation}

and summing the correlation over a complete $\qvec$-grid of $\Na$ points,

\begin{equation}
    \sum_{\qvec} s^\alpha(\qvec,t{+}\tau)\,s^{\beta*}(\qvec,t)
    = \frac{1}{\Na}\sum_{\qvec}\sum_{jk}
        \field_j^\alpha(t{+}\tau)\field_k^\beta(t)\,e^{i\qvec\cdot(\rvec_j-\rvec_k)} ,
\end{equation}

the $\qvec$-sum collapses through
$\sum_{\qvec}e^{i\qvec\cdot(\rvec_j-\rvec_k)} = \Na\,\delta_{jk}$, leaving

\begin{equation}
\boxed{
    \sum_{\qvec} s^\alpha(\qvec,t{+}\tau)\,s^{\beta*}(\qvec,t)
    = \sum_j \field_j^\alpha(t{+}\tau)\,\field_j^\beta(t).
}
\label{eq:parseval}
\end{equation}

The Brillouin-zone integral of the correlation is exactly the sum of the
\emph{on-site} correlations. No $\qvec$ appears on the right-hand side, no
$\qvec$-mesh has to be chosen, and nothing is approximated: this is Parseval's
identity, not a sampling scheme.

The chain is therefore

\begin{equation}
    \field_j^\alpha(t) \;\longrightarrow\; x_j(t)
    \;\longrightarrow\; C(\tau) \;\longrightarrow\; D(\omega),
\end{equation}

where $x_j$ is a scalar time series per atom, obtained by projecting the field
onto whichever direction in spin space the requested channel names
(Sec.~\ref{sec:channels}).

\begin{remark}[Cost]
The $\qvec$-mesh route costs $O(N_q\,\Nt\,\Na)$ for the spatial transforms
alone. This route costs $O(\Na\,\Nt\log\Nt)$ in total --- one transform per
atom, and no $\qvec$-loop at all.
\end{remark}

\begin{remark}[Code layout]
The four lines of \texttt{dos\_for\_channel} in \texttt{sqw\_dos.py} are the
whole calculation. The numerics live in \texttt{dos.py}; the channel model is
shared with \texttt{channels.py} and the window with \texttt{correlate.py}, so
the conventions are identical to those of \texttt{sqw\_spin\_corr.py} (see
\texttt{doc\_sqw\_spin\_corr.tex}). There is no MPI: the whole calculation is
cheaper than reading the dump.
\end{remark}

% ======================================================================
\section{Input}
% ======================================================================

The same trajectory as the $\qvec$-resolved route: a three-component field
$\field_j^\alpha(t)$ on $\Na$ atoms over $\Nt$ frames, optionally mass-weighted
by $w_j = \sqrt{m_{\text{type}(j)}}$ and optionally filtered by
$\max_t\lVert\field_j\rVert \le \tau_{\text{spin}}$.

Atomic \emph{positions} are never read. Neither is the reciprocal lattice ---
there is no $\qvec$ to express --- so no supercell need be supplied.

% ======================================================================
\section{Channels}
\label{sec:channels}
% ======================================================================

Which component of the spin should be correlated? That is the only choice to
make, and it is expressed exactly as in the $\qvec$-resolved route: a channel
is a real $3\times3$ weight matrix $\Wmat$, and the quantity it names is

\begin{equation}
\boxed{
    \sum_{ab} W_{ab}\,C^{ab}(\tau),
    \qquad
    C^{ab}(\tau) = \sum_j \bigl\langle \field_j^a(t{+}\tau)\,\field_j^b(t)\bigr\rangle .
}
\label{eq:channel}
\end{equation}

The \texttt{-\/-component} grammar is unchanged: tokens separated by spaces
produce separate curves, terms joined by \texttt{+} are summed into one, and
the default \texttt{1+5+9} is the trace.

\subsection{Two Restrictions}

\textbf{\texttt{L} and \texttt{T} are unavailable.} Their weights are
$\qhat\qhat^{\mathsf T}$ and $\Imat-\qhat\qhat^{\mathsf T}$, built from the
\emph{direction} of $\qvec$. The DOS has integrated $\qvec$ away, so there is
no direction left to project on. Requesting them raises an error rather than
silently returning something else.

\textbf{$\Wmat$ must be symmetric.} A non-symmetric weight --- a bare
off-diagonal element such as \texttt{2}, or a mixed group such as \texttt{1+2}
--- names a signed quantity, not a spectral density. A density of states is
non-negative by definition, so these are rejected. This is the same
positive-semidefiniteness criterion that decides clipping in the
$\qvec$-resolved route, applied one step earlier.

\subsection{Evaluation by Eigen-Decomposition}

Equation~\eqref{eq:channel} is not evaluated by building the nine
cross-correlations and contracting them. Because $\Wmat$ is symmetric it can be
diagonalised, $\Wmat = \sum_i \lambda_i \vvec_i\vvec_i^{\mathsf T}$, and then

\begin{align}
    \sum_{ab} W_{ab}\,C^{ab}(\tau)
    &= \sum_j \sum_i \lambda_i \sum_{ab} v_i^a v_i^b
        \bigl\langle \field_j^a(t{+}\tau)\field_j^b(t)\bigr\rangle \notag\\
    &= \sum_i \lambda_i \sum_j
        \bigl\langle x_j^{(i)}(t{+}\tau)\,x_j^{(i)}(t)\bigr\rangle,
    \qquad x_j^{(i)} = \vvec_i\cdot\bm{\field}_j .
    \label{eq:eig}
\end{align}

Every term is now an ordinary autocorrelation of a real \emph{scalar} series.
Two consequences:

\begin{itemize}
    \item only $\operatorname{rank}(\Wmat)$ transforms are needed, at most
    three, instead of up to nine cross-correlations;
    \item exactly one spectrum is resident at a time, so the memory high-water
    mark is one $(\!\sim\!\Nt,\Na)$ array rather than two or more. At
    $\Nt = 10^4$ and $\Na = 10^3$ that is the difference between roughly
    $260$ MiB and $780$ MiB.
\end{itemize}

For the default trace $\Wmat = \Imat$ the eigenvectors are the Cartesian axes
and Eq.~\eqref{eq:eig} reduces to the obvious procedure: autocorrelate
$\field^x$, $\field^y$, $\field^z$ and add.

\begin{center}
\begin{tabular}{lllc}
    \toprule
    Token & $\Wmat$ & Directions $\vvec_i$ & Transforms \\
    \midrule
    \texttt{1+5+9} (default) & $\Imat$ & $\hat x,\hat y,\hat z$ & 3 \\
    \texttt{1} & $\operatorname{diag}(1,0,0)$ & $\hat x$ & 1 \\
    \texttt{1+5} & $\operatorname{diag}(1,1,0)$ & $\hat x,\hat y$ & 2 \\
    \texttt{2+4} & $e_xe_y^{\mathsf T}+e_ye_x^{\mathsf T}$ & $(\hat x\pm\hat y)/\sqrt2$, $\lambda=\pm1$ & 2 \\
    \bottomrule
\end{tabular}
\end{center}

The last row shows that the restriction to symmetric $\Wmat$ is not a
restriction to diagonal ones: a symmetric off-diagonal combination is perfectly
computable, it simply corresponds to projecting along a rotated pair of axes.

% ======================================================================
\section{Preparing the Series}
% ======================================================================

Before correlating, each projected series $x_j(t)$ has its time average removed
and is then tapered:

\begin{equation}
    \tilde x_j(t) = w(t)\,\bigl[x_j(t) - \langle x_j\rangle_t\bigr],
    \qquad t = 0,\dots,\Nt-1 .
\end{equation}

The order matters: tapering a series that still carries a DC offset would
imprint the shape of the window onto the spectrum.

$w$ is a \textbf{data-domain} window --- Hann by default, or unity with
\texttt{-\/-window none} --- applied to the time series, not to the
correlation. Its purpose is to suppress spectral leakage: the record is finite
and the transform treats it as periodic, so unless a mode completes a whole
number of cycles the splice between last and first frame is a discontinuity
whose power smears across all frequencies.

\begin{remark}[Why no lag window]
Tapering $C(\tau)$ instead would reduce the statistical scatter rather than the
leakage. That is unnecessary here for two reasons. The DOS already sums over
$\Na$ sites, so its scatter is smaller than that of a single-$\qvec$ spectrum
by roughly $\sqrt{\Na}$ --- for $10^3$ atoms the curve is already smooth. And
where further smoothing is wanted, \texttt{-\/-smooth-sigma-thz} provides it
explicitly and visibly, with a width the user controls, instead of an implicit
taper. Adding a lag window on top would smooth twice.
\end{remark}

% ======================================================================
\section{The Correlation}
% ======================================================================

\begin{equation}
\boxed{
    C(\tau) = \frac{1}{M_\tau\,\Na}\sum_i \lambda_i
              \sum_j \sum_t \tilde x_j^{(i)}(t{+}\tau)\,\tilde x_j^{(i)}(t),
    \qquad \tau = 0,\dots,\Nt-1,
}
\label{eq:corr}
\end{equation}

with $M_\tau = \Nt$ for \texttt{-\/-corr-norm biased} (the default, which keeps
the estimate a genuine power spectrum) or $M_\tau = \Nt-\tau$ for
\texttt{unbiased} (which removes the triangular bias at the cost of positive
semidefiniteness).

\subsection{Evaluation}

By the Wiener--Khinchin identity, with the series zero-padded to
$L \ge 2\Nt$,

\begin{equation}
\boxed{
    \sum_t \tilde x(t{+}\tau)\tilde x(t)
    = \mathcal{F}^{-1}\bigl[\,|\hat x|^2\,\bigr](\tau).
}
\end{equation}

The padding makes the circular correlation of the padded signal equal the
linear correlation of the original, so this is exact. The code takes
$L = 2^{\lceil \log_2(2\Nt-1)\rceil}$.

The projected field is real, so the half-spectrum routines
\texttt{rfft}/\texttt{irfft} do the work at half the cost of a complex pair.
Everything runs in float64 regardless of the storage dtype of the trajectory:
the sum runs over $\Nt$ terms and would otherwise eat into float32's seven
digits.

\subsection{Normalisation Convention}

The $1/\Na$ in Eq.~\eqref{eq:corr} is \textbf{per atom}, not per
atom-times-component. This matters when comparing channels: with a per-atom
normalisation the identity

\begin{equation}
    D_{\texttt{1+5+9}}(\omega)
    = D_{\texttt{1}}(\omega) + D_{\texttt{5}}(\omega) + D_{\texttt{9}}(\omega)
\end{equation}

holds exactly, whereas dividing by the number of contributing components would
scale the trace down by three relative to the sum of its parts and make
absolute values incomparable between channels. With
\texttt{-\/-normalize max} or \texttt{area} the constant is absorbed anyway;
it is visible only under \texttt{-\/-normalize none}.

% ======================================================================
\section{Frequency Domain}
% ======================================================================

$C(\tau)$ is real and even --- projecting a real field on a real direction
leaves nothing to conjugate --- so the extension to negative lags is a plain
mirror, and

\begin{equation}
\boxed{
    D(\omega_k) = \sum_{\tau=-(\Nt-1)}^{\Nt-1}
                  e^{-i\omega_k\tau\Delta t}\,C(|\tau|),
    \qquad \omega_k \ge 0,
}
\end{equation}

evaluated by FFT after \texttt{ifftshift} maps the two-sided sequence into FFT
order. There is no $\Delta t$ prefactor, so $D$ is a relative intensity. The
frequency grid follows the transform length $L = 2\Nt-1$,

\begin{equation}
    f_k = \frac{k}{(2\Nt-1)\,\Delta t\,[\mathrm{s}]\times10^{12}}\ \mathrm{THz},
    \qquad k = 0,\dots,\Nt-1 ,
\end{equation}

and negative values are clipped to zero when the channel weight is positive
semidefinite, which is the exact condition for the result to be a physical
intensity.

% ======================================================================
\section{Presentation}
% ======================================================================

Three optional steps are applied to the raw curve, in this order:

\begin{enumerate}
    \item \textbf{Trim} to $[f_{\min}, f_{\max}]$
    (\texttt{-\/-freq-min-thz}, \texttt{-\/-freq-max-thz}).
    \item \textbf{Smooth} with a Gaussian of width
    \texttt{-\/-smooth-sigma-thz}, converted to a number of grid points using
    the median spacing of the trimmed grid. Zero disables it.
    \item \textbf{Normalise} (\texttt{-\/-normalize}): \texttt{max} divides by
    the peak, \texttt{area} by $\int D\,df$, \texttt{none} leaves the raw
    scale.
\end{enumerate}

Both the raw and the finished curve are stored, so the effect of steps 2--3 can
always be inspected; \texttt{-\/-plot-raw} overlays them.

% ======================================================================
\section{Validation}
% ======================================================================

Checked on \texttt{examples/s4.lammpstrj}, a synthetic trajectory whose spins
precess coherently at $4.0$~THz in the $xy$ plane ($2048$ frames, $400$ atoms,
$\Delta t = 2$~fs):

\begin{center}
\begin{tabular}{lll}
    \toprule
    Check & Result & \\
    \midrule
    Trace $=$ sum of diagonal channels
        & $3.5\times10^{-16}$ & machine precision \\
    Peak of \texttt{1+5+9}, \texttt{1}, \texttt{5}
        & $4.0293$ THz & the injected mode \\
    Peak of \texttt{9} ($\field^{zz}$)
        & $0$ THz & spins precess in $xy$; $\field^z$ is static \\
    Against the previous implementation
        & $2.0\times10^{-10}$ & after removing the factor $3$ \\
    \bottomrule
\end{tabular}
\end{center}

The third row is the useful one: it confirms that the channel selection really
does isolate a Cartesian component, since a mode confined to the $xy$ plane
must leave $\field^{zz}$ with no dynamics at all.

% ======================================================================
\section{Summary of the Algorithm}
% ======================================================================

Every step below rests on Eq.~\eqref{eq:parseval}: the Brillouin-zone sum has
already been done analytically, so no $\qvec$ ever appears.

\begin{enumerate}
    \item \textbf{Parse channels} --- reject \texttt{L}/\texttt{T} and any
    non-symmetric weight.
    \item \textbf{Diagonalise} each $\Wmat$ into directions $\vvec_i$ with
    weights $\lambda_i$, Eq.~\eqref{eq:eig}.
    \item For each direction: \textbf{project} $x_j = \vvec_i\cdot\bm{\field}_j$,
    \textbf{subtract the mean}, \textbf{apply the data window}.
    \item \textbf{Autocorrelate} by zero-padded \texttt{rfft}/\texttt{irfft},
    sum over atoms, accumulate with weight $\lambda_i$.
    \item \textbf{Normalise} by $M_\tau \Na$.
    \item \textbf{Mirror} to negative lags and \textbf{transform}; clip if
    $\Wmat$ is positive semidefinite.
    \item \textbf{Trim, smooth, normalise} for presentation.
\end{enumerate}

% ======================================================================
\section{Notation}
% ======================================================================

\begin{table}[h]
    \centering
    \begin{tabular}{ll}
        \toprule
        Symbol & Meaning \\
        \midrule
        $\Na$ & Atoms retained after the field threshold \\
        $\Nt$ & Time frames \\
        $\Delta t$ & Interval between saved frames (fs) \\
        $\field_j^\alpha(t)$ & Field component $\alpha$ on atom $j$ \\
        $\Wmat$ & Channel weight matrix, Eq.~\eqref{eq:channel} \\
        $\lambda_i,\ \vvec_i$ & Eigenvalues and directions of $\Wmat$, Eq.~\eqref{eq:eig} \\
        $x_j^{(i)}$ & Projected scalar series $\vvec_i\cdot\bm{\field}_j$ \\
        $w(t)$ & Data window \\
        $C(\tau)$ & On-site correlation of the channel, Eq.~\eqref{eq:corr} \\
        $M_\tau$ & Correlation denominator: $\Nt$ or $\Nt-\tau$ \\
        $D(\omega)$ & Density of states \\
        \bottomrule
    \end{tabular}
    \caption{Notation.}
\end{table}

\end{document}
